\documentclass[%
a4paper,
DIV12,
2.5headlines,
bigheadings,
titlepage,
openbib,
%draft
]{scrartcl}

\usepackage{amsfonts}
\usepackage{amsmath}
%\usepackage{lm}
\usepackage{ifxetex}
\usepackage{ifluatex}
\usepackage{eurosym}
\usepackage{listings}
\usepackage{fancyvrb}
\usepackage{longtable}
\usepackage{booktabs}
\usepackage{graphicx}
\usepackage{grffile}
\usepackage{ulem}
\usepackage{tabularx}
\usepackage[export]{adjustbox}
\usepackage[T1]{fontenc}
\usepackage{listings}
\usepackage{titling}
\usepackage{graphicx}
\usepackage{scrpage2}
\usepackage{float}
\usepackage[colorinlistoftodos,prependcaption,textsize=tiny]{todonotes}
\usepackage[pdftex, colorlinks, linktocpage, linkcolor=black, citecolor=black, urlcolor=black]{hyperref}
\usepackage{atbegshi}% http://ctan.org/pkg/atbegshi
\usepackage{placeins}
\usepackage[ampersand]{easylist}

\AtBeginDocument{\AtBeginShipoutNext{\AtBeginShipoutDiscard}}
\pagestyle{scrheadings}

\pretitle{%
  \begin{center}
  \LARGE
  \includegraphics[width=6cm]{figures/octocat.jpg}\\[\bigskipamount]
}
\posttitle{\end{center}}

\begin{document}

\providecommand{\tightlist}{%
  \setlength{\itemsep}{0pt}\setlength{\parskip}{0pt}}

\DeclareRobustCommand{\desiredTitle}{
  Intellitank - Prediction of Gas Prices and Route Optimization \\\normalsize{Solution for the informatiCup 2018}
}

\DeclareRobustCommand{\desiredAuthor}{
  \\\\Freya Behrens\\\texttt{freya.behrens@student.hpi.de}\\Hasso-Plattner-Institute \and \\\\Sebastian Bischoff\\\texttt{sebastian.bischoff@student.hpi.de}\\Hasso-Plattner-Institute \and \\\\Willi Gierke\\\texttt{willi.gierke@student.hpi.de}\\Hasso-Plattner-Institute\\\\\\\\\\\\
}

\begin{titlepage}
\begin{center}
   \title{\desiredTitle}
   \author{\desiredAuthor}
   \date{Potsdam, January 2018}
   \maketitle
\end{center}
\end{titlepage}

\pagebreak

\tableofcontents
\pagebreak

\section{Challenge Description}\label{challenge-description}
\begin{itemize}
\item Describe task
\end{itemize}

\section{Data Exploration}\label{data-exploration}
\begin{itemize}
\item Some plots of data\\
\item Characteristic plots (with outliers etc.)
\item Calculate correlation between gas price and different crude oil and gasoline markets (gas price as average/median over all German gas stations or average correlation over all gas stations)
\end{itemize}


\subsection{Data Retrieval}\label{data-retrieval}
Additional data:
\begin{itemize}
\item clustering

\begin{itemize}
\item geolocation
\item owning company
\item price change behavior
\end{itemize}

\item world market

\begin{itemize}
\item barrel oil
\item gas?
\end{itemize}

\item political features

\begin{itemize}
\item pipeline opening/closing
\item sanctions
\item OPEC
\end{itemize}

\item opening of new gas station
\item closing of old gas station
\item very close gas stations
 
\begin{itemize}
\item on other Autobahn lane
\item on opposite sides of the same street
\end{itemize}

\item temperature
\item weather
\item news (?)
\item holidays
\item price in neighboring country
\item on traveling route to other countries
\item distance to Autobahn
\item traffic information
\item traffic density
\end{itemize}
 


\subsection{Data Analysis}\label{data-analysis}



\section{Prediction Model}\label{prediction-model}
\begin{itemize}
\item Maybe first predict time and then delta of change\\
\item \textbf{Baseline}: Predict time series only with historical data without further data (facebook prophet)\\

\begin{itemize}
\item Doesn't take interactions between gas stations into account because one learns for each gas station a separate model
\end{itemize}
		
\item How to incorporate static data like distance to Autobahn? Cluster similar gas stations and learn joint prediction model. 
What is the advantage over single models? Puhhh...\\
\item Big ensemble:

\begin{itemize}
\item One global model which predicts average gas price (change)
\item One cluster model which predicts change in cluster of similar gas stations relative to global model

\begin{itemize}
\item Cluster based on static attributes

\begin{itemize}
\item Assumption would be that the behavior of the gas price is dependent on the static attributes --> There is a causal connection between static attributes and behavior of gas price
\item e.g. all gas stations on Autobahnen share the same behavior
\end{itemize}

\item Cluster based on static attributes and single value features of a time series (e.g. average number of price changes per day, time of maximum price etc.)

\begin{itemize}
\item Doesn't have clear theoretical explanation yet
\end{itemize}

\item Cluster based on time series 

\begin{itemize}
\item Time series are only described by the behavior of the time series and not by static attributes
\item e.g. only dependent on risk disposition of gas station attendant and not on vicinity of Autobahn
\item Could cluster for example based on correlation between time series
\item look into literature
\item ask someone at the chair
\item No use for static features?
\end{itemize}
\end{itemize}

\item One gas station model which predicts change of individual gas station relative to cluster model
\end{itemize}
\end{itemize}



\subsection{Training and Test Data
Set}\label{training-and-test-data-set}


\subsection{Data Cleaning and
Preprocessing}\label{data-cleaning-and-preprocessing}


\subsection{Feature Generation from Existing
Data}\label{feature-generation-from-existing-data}


\subsection{Feature Selection}\label{feature-selection}



\subsection{Feature Engineering}\label{feature-engineering}



\subsection{Validation of Prediction
Model}\label{validation-of-prediction-model}
\begin{itemize}
\item Evaluate on fixed time span? Like last month?
\item Data doesn't have discrete change intervals like each minute, hour, day...
\item Need fitting cost function

\begin{itemize}
\item Normal MSE
\item MSE but algorithm is restricted to maximum $n$ price changes per day
\end{itemize}

\item One could train it as "end-to-end-model"

\begin{itemize}
\item Evaluate on multiple routes in conjunction with defined times for refueling stops
\item Routes could be generated with random walks on gas station graph
\item Optimization could be challenging
\end{itemize}
\end{itemize}


\subsection{Overall Prediction Model}\label{overall-prediction-model}

\section{Algorithm for refueling}
"To fill or not to fill"\cite{khuller2007fill}\\
For extension to electric cars see \ref{extensions}

\section{Implemented Application}\label{implemented-application}
\subsection{Tooling Choices}\label{tooling-choices}



\subsection{Functionality}\label{functionality}



\section{Validation}\label{validation}
\begin{itemize}
\item Discuss results of prediction (up to one month ahead)
\item Evaluate with suitable metrics
\end{itemize}


\section{Extensions}\label{extensions}
\begin{itemize}
\item App where you can specify a route, it calculates the optimal refuel stops and visualizes it on a map
\item Model the problem for electric cars

\begin{itemize}
\item we can assume that the price per kWh is equal at each charger (check this)\\
\item but different charger charge with a different speed (2 to 22(?)kW)\\
\item not every charger is with every car compatible\\
\item prediction task: predict if the charger is occupied or free\\
\end{itemize}

\item Map or list where you can search gas stations in your vicinity and check if they are more expensive/cheaper than other gas stations
\end{itemize}

\bibliographystyle{acm}
\bibliography{documentation} 

\end{document}
